\documentclass[A4paper,11pt]{article}
\usepackage{graphics,graphicx,epsfig,wrapfig}
\usepackage{longtable}  
\usepackage{amsmath,verbatim,amssymb,color,lscape}
\usepackage{kotex}
\usepackage{natbib}
\usepackage{lastpage}
\usepackage{multirow} 
\usepackage{authblk}
\usepackage{bm}
\usepackage{subcaption}
\usepackage{algorithm}
\usepackage{algpseudocode}
\usepackage{indentfirst}
\usepackage[ruled,vlined]{algorithm2e}
\usepackage{algorithmic}
\usepackage[normalem]{ulem}
\usepackage[export]{adjustbox}
\usepackage{tikz}   
\usepackage{float}
\usepackage{listings}
\usepackage{color}
\usepackage{subcaption}  % 프리앰블에 추가
\definecolor{codegreen}{rgb}{0,0.6,0} % 프로그래밍시 볼 수 있는 그 초록색입니다.

% Define a custom style
\lstdefinestyle{mystyle}{
    backgroundcolor=\color{white}, % 배경색 지정 (흰색 선호) 
    commentstyle=\color{codegreen}, % 주석의 글자 색 지정
    keywordstyle=\color{blue}, % 키워드의 글자 색 지정
    stringstyle=\color{codegreen}, % 문자열의 글자 색
    basicstyle=\ttfamily\footnotesize, % 글자 스타일 (작고 굵어짐)
    breaklines=true, % 자동 line breaking
    captionpos=b, % 코드 블럭의 제목 위치
    numbers=left, % 코드 줄 넘버 표시 (왼쪽)
    keepspaces=true }
\lstset{style=mystyle}
\usepackage{pgffor}         % \foreach 반복문

%\usepackage{hyperref}
\usepackage{verbatim}
\newcommand{\hide}[1]{}  
\newcommand{\msim}{\mathop{\rm \sim}}
\newcommand{\ind}{\msim\limits^{\mbox{\tiny ind}}}
\newcommand{\iid}{\msim\limits^{\mbox{\tiny iid}}}
\makeatletter
\newcommand{\s}{\vspace{.25cm}}
\renewcommand{\familydefault}{\sfdefault}
\usepackage[outerbars,color]{changebar}
\ifx\pdfoutput\undefined
\else\ifnum\pdfoutput>0
  \usepackage{pdfcolmk}
\fi\fi
\cbcolor{black}

\setlength{\changebarsep}{5mm}

\usepackage[margin=1.0in]{geometry}
\baselineskip=15.5pt

\usepackage{tikz}
\usetikzlibrary{bayesnet}
\usetikzlibrary{fit,positioning}
\newcommand{\obs}{latent}
\newcommand{\lat}{obs}
\newcommand{\con}{const}
\newcommand{\independent}{\perp\mkern-9.5mu\perp}

\setcounter{secnumdepth}{4}
\renewcommand{\baselinestretch}{1.5}

\newfont{\rmm}{cmr10 at 11pt}
\rmm
\newcommand{\widesim}[2][1.5]{
  \mathrel{\overset{#2}{\scalebox{#1}[1]{$\sim$}}}
}
%\appendix

\pagestyle{plain}

\title{bipartite SBM}
\author[]{윤현석}
\date{\today}

\begin{document}
\maketitle
\section{bipartite SBM}

% =========================================================
% Methodology section
% =========================================================

\subsection{Flat weighted bipartite SBM for strictly positive distance-type edges}

본 절에서는 strictly positive edge weight를 가지는 bipartite network를 위한 flat weighted stochastic block model (SBM)을 제안한다. 관측 데이터는 $n \times p$ 행렬
\[
X = (X_{ij}) \in \mathbb{R}_{+}^{n\times p},
\qquad X_{ij} > 0,
\]
로 주어지며, 각 원소 $X_{ij}$는 row node $i=1,\ldots,n$와 column node $j=1,\ldots,p$ 사이의 거리형(distance-type) edge weight를 의미한다. 본 연구의 데이터 구조에서는 edge weight의 절댓값이 작을수록 두 노드가 더 유사하다고 해석한다. 본 모형은 row-side partition과 column-side partition을 동시에 추정하는 co-clustering 모형이지만, 실질적인 관심의 초점은 column-side clustering에 있으므로 최종 해석은 주로 $\{w_j\}_{j=1}^p$를 중심으로 수행한다.

\subsubsection{Model structure}

row node $i$와 column node $j$의 잠재 군집 할당을 각각
\[
z_i \in \{1,\ldots,Q\},
\qquad
w_j \in \{1,\ldots,L\}
\]
로 둔다. 여기서 $Q$는 row cluster의 개수, $L$은 column cluster의 개수이다. 따라서 잠재 membership 벡터는
\[
\mathbf{z} = (z_1,\ldots,z_n) \in \{1,\ldots,Q\}^{n},
\qquad
\mathbf{w} = (w_1,\ldots,w_p) \in \{1,\ldots,L\}^{p}
\]
이다.

각 block pair $(q,l)$에 대해, strictly positive edge weight를 모형화하기 위하여 Gamma likelihood를 가정한다:
\[
X_{ij} \mid (z_i=q,\; w_j=l,\; \kappa,\; \lambda_{ql})
\sim \operatorname{Gamma}(\kappa,\lambda_{ql}),
\qquad q=1,\ldots,Q,\quad l=1,\ldots,L,
\]
여기서 $\kappa>0$는 global shape parameter이고, $\lambda_{ql}>0$는 block-specific rate parameter이다. 이에 대응하는 density는
\[
f(X_{ij}\mid z_i=q,w_j=l,\kappa,\lambda_{ql})
=
\frac{\lambda_{ql}^{\kappa}}{\Gamma(\kappa)}
X_{ij}^{\kappa-1}
\exp(-\lambda_{ql}X_{ij}),
\qquad X_{ij}>0,
\]
이다.

이때 block $(q,l)$의 조건부 평균은
\[
\mu_{ql}
=
E(X_{ij}\mid z_i=q,w_j=l,\kappa,\lambda_{ql})
=
\frac{\kappa}{\lambda_{ql}}
\]
이므로,
\[
\lambda_{ql}\uparrow
\quad\Longleftrightarrow\quad
\mu_{ql}\downarrow
\]
이다. 즉, $\lambda_{ql}$가 클수록 해당 row cluster $q$와 column cluster $l$ 사이의 평균 거리가 작아지므로 두 cluster 간 유사성이 높은 것으로 해석된다.

본 모형의 주요 모수는 다음과 같이 정리할 수 있다:
\[
\pi = (\pi_1,\ldots,\pi_Q) \in \Delta^{Q-1},
\qquad
\rho = (\rho_1,\ldots,\rho_L) \in \Delta^{L-1},
\]
\[
\Lambda = (\lambda_{ql})_{q=1,\ldots,Q;\,l=1,\ldots,L}\in\mathbb{R}_{+}^{Q\times L},
\qquad
\kappa \in \mathbb{R}_{+},
\]
where $\Delta^{K-1}$ denotes the $(K-1)$-dimensional probability simplex.

조건부 독립 가정 하에서 complete-data likelihood는
\[
p(X\mid \mathbf{z},\mathbf{w},\kappa,\Lambda)
=
\prod_{i=1}^{n}\prod_{j=1}^{p}
f(X_{ij}\mid z_i,w_j,\kappa,\Lambda)
=
\prod_{i=1}^{n}\prod_{j=1}^{p}
\frac{\lambda_{z_iw_j}^{\kappa}}{\Gamma(\kappa)}
X_{ij}^{\kappa-1}
\exp(-\lambda_{z_iw_j}X_{ij})
\]
로 쓸 수 있다.

\subsubsection{Prior specification}

cluster proportion에 대해서는 대칭 Dirichlet prior를 사용한다:
\[
\pi \sim \operatorname{Dirichlet}(1,\ldots,1),
\qquad
\rho \sim \operatorname{Dirichlet}(1,\ldots,1).
\]
그리고 latent membership는
\[
z_i\mid \pi \sim \operatorname{Categorical}(\pi),
\qquad i=1,\ldots,n,
\]
\[
w_j\mid \rho \sim \operatorname{Categorical}(\rho),
\qquad j=1,\ldots,p,
\]
로 둔다.

block-specific rate parameter에 대해서는 데이터 전체 평균
\[
\bar X = \frac{1}{np}\sum_{i=1}^{n}\sum_{j=1}^{p}X_{ij}
\]
를 이용하여 다음의 Gamma prior를 사용한다:
\[
\lambda_{ql}\mid \kappa
\sim
\operatorname{Gamma}(\kappa r,\; r\bar X),
\qquad q=1,\ldots,Q,\quad l=1,\ldots,L,
\]
여기서 $r=0.1$은 prior strength이다. 이 설정은 각 block에 대해 약 $0.1$개의 pseudo-edge에 해당하는 약한 정규화 효과를 부여하며, prior mean은
\[
E(\lambda_{ql}\mid \kappa)=\frac{\kappa r}{r\bar X}=\frac{\kappa}{\bar X}
\]
이 된다. 따라서 prior는 각 block의 평균 거리 수준이 전체 평균 거리와 크게 다르지 않다는 약한 정보만 제공한다.

global shape parameter는 양수 제약을 반영하기 위하여 log scale에서 정규 prior를 둔다:
\[
\log \kappa \sim \mathcal{N}(\mu_{\log\kappa},\sigma_{\log\kappa}^{2}),
\qquad
\mu_{\log\kappa}=0,\quad
\sigma_{\log\kappa}^{2}=4.
\]
즉, $\kappa$의 prior median은 1이며, 충분히 넓은 범위를 허용하는 weakly informative prior이다.

이상의 prior를 종합하면 joint distribution은
\begin{align*}
p(X,\mathbf{z},\mathbf{w},\pi,\rho,\Lambda,\kappa)
&=
p(X\mid \mathbf{z},\mathbf{w},\Lambda,\kappa)\,
p(\mathbf{z}\mid\pi)\,
p(\mathbf{w}\mid\rho)\,
p(\pi)\,
p(\rho)\,
p(\Lambda\mid \kappa)\,
p(\kappa) \\
&=
\left[
\prod_{i=1}^{n}\prod_{j=1}^{p}
f(X_{ij}\mid z_i,w_j,\kappa,\Lambda)
\right]
\left[
\prod_{i=1}^{n}\pi_{z_i}
\right]
\left[
\prod_{j=1}^{p}\rho_{w_j}
\right] \\
&\quad\times
\operatorname{Dir}(\pi;1,\ldots,1)\,
\operatorname{Dir}(\rho;1,\ldots,1)\,
\prod_{q=1}^{Q}\prod_{l=1}^{L}
p(\lambda_{ql}\mid \kappa)\,
p(\kappa).
\end{align*}

\subsubsection{Posterior inference}

사후분포는 다음과 같이 주어진다:
\[
p(\mathbf{z},\mathbf{w},\pi,\rho,\Lambda,\kappa \mid X)
\propto
p(X\mid \mathbf{z},\mathbf{w},\Lambda,\kappa)\,
p(\mathbf{z}\mid\pi)\,
p(\mathbf{w}\mid\rho)\,
p(\pi)\,
p(\rho)\,
p(\Lambda\mid \kappa)\,
p(\kappa).
\]

본 모형에서는 일부 모수의 완전 조건부 분포가 닫힌 형태를 가지므로 Gibbs sampling과 Metropolis--Hastings (MH) 업데이트를 결합한 MCMC 알고리즘을 사용한다. 구체적으로, $(\pi,\rho,\Lambda)$는 Gibbs 단계로 갱신하고, $(\mathbf{z},\mathbf{w})$는 categorical Gibbs 단계로 갱신하며, $\kappa$는 log scale random-walk MH로 갱신한다.

실질적 관심이 column-side clustering에 있으므로, posterior summary에서는 저장된 MCMC 표본으로부터 column co-clustering probability
\[
\widehat{P}(w_j = w_{j'} \mid X)
\]
를 계산하고, 이를 기반으로 최종 column partition을 요약한다. row-side clustering은 보조적 구조로 함께 추정되지만, 해석의 초점은 $\mathbf{w}$와 $\Lambda$가 정의하는 block structure에 둔다.

\subsubsection{Identifiability and label switching}

본 모형은 latent label의 순열(permutation)에 대하여 불변이다. 즉, row cluster label과 column cluster label은 각각 독립적인 재명명(relabeling)에 대해 동일한 likelihood를 생성하므로, MCMC 표본에는 label switching 현상이 발생할 수 있다. 본 연구에서는 이를 사전적 제약으로 처리하지 않고, MCMC 종료 후 사후적 relabeling 절차를 적용한다.

특히 본 연구의 주된 관심이 column-side partition에 있으므로, 저장된 각 반복에서 먼저 column label $\mathbf{w}$를 기준 참조 partition에 정렬한 다음, 이에 조건부하여 row label $\mathbf{z}$를 정렬한다. 이때 block-specific parameter matrix $\Lambda$ 역시 동일한 row/column permutation을 적용하여 함께 재배열한다. 최종 posterior summary와 block interpretation은 이러한 post-processing을 거친 표본을 기준으로 계산한다.

\subsubsection{Model selection via ICL}

cluster 수 $(Q,L)$는 integrated completed likelihood (ICL) criterion에 기반하여 선택한다. 후보 집합
\[
Q \in \mathcal{Q}=\{1,\ldots,Q_{\max}\},
\qquad
L \in \mathcal{L}=\{1,\ldots,L_{\max}\}
\]
을 설정하고, 각 $(Q,L)$에 대하여 모형을 적합한 뒤 다음의 ICL 값을 계산한다:
\[
\operatorname{ICL}(Q,L)
=
\log p(X,\hat{\mathbf{z}},\hat{\mathbf{w}} \mid \hat{\Theta}_{Q,L},Q,L)
-
\frac{\nu_{Q,L}}{2}\log(np),
\]
where
\[
\hat{\Theta}_{Q,L}=(\hat{\pi},\hat{\rho},\hat{\Lambda},\hat{\kappa}),
\qquad
\nu_{Q,L}=(Q-1)+(L-1)+QL+1.
\]
여기서 $\hat{\mathbf{z}},\hat{\mathbf{w}}$는 적합된 모형에서의 대표 partition이며, $\nu_{Q,L}$는 자유도에 해당한다. complete-data log-likelihood term은
\begin{align*}
\log p(X,\hat{\mathbf{z}},\hat{\mathbf{w}} \mid \hat{\Theta}_{Q,L},Q,L)
&=
\sum_{i=1}^{n}\log \hat{\pi}_{\hat z_i}
+
\sum_{j=1}^{p}\log \hat{\rho}_{\hat w_j} \\
&\quad+
\sum_{i=1}^{n}\sum_{j=1}^{p}
\log f(X_{ij}\mid \hat z_i,\hat w_j,\hat\kappa,\hat\Lambda).
\end{align*}
최종적으로
\[
(\hat Q,\hat L)
=
\arg\max_{(Q,L)\in\mathcal{Q}\times\mathcal{L}}
\operatorname{ICL}(Q,L)
\]
을 선택한다. 본 연구에서는 column-side clustering이 주요 관심 대상이므로, 실제 탐색에서는 $L$에 대해 $Q$보다 더 촘촘한 grid를 두는 전략을 사용할 수 있다.

% =========================================================
% Appendix section
% =========================================================

\section{appendix}

\subsection{model description}

\subsubsection{prior distribution}

사전 분포(prior distribution)는 다음과 같이 설정한다.

\begin{table}[H]
\centering
\begin{tabular}{c|l|l}
\hline
\textbf{번호} & \textbf{분포} & \textbf{인덱스} \\ \hline

1 &
$\displaystyle \pi \sim \operatorname{Dirichlet}(1,\ldots,1)$ &
$\pi \in \Delta^{Q-1}$ \\ \hline

2 &
$\displaystyle \rho \sim \operatorname{Dirichlet}(1,\ldots,1)$ &
$\rho \in \Delta^{L-1}$ \\ \hline

3 &
$\displaystyle z_i \mid \pi \overset{\text{ind}}{\sim} \operatorname{Categorical}(\pi)$ &
$i=1,\ldots,n$ \\ \hline

4 &
$\displaystyle w_j \mid \rho \overset{\text{ind}}{\sim} \operatorname{Categorical}(\rho)$ &
$j=1,\ldots,p$ \\ \hline

5 &
$\displaystyle \lambda_{ql}\mid \kappa \overset{\text{ind}}{\sim}
\operatorname{Gamma}(\kappa r,\; r\bar X)$ &
$q=1,\ldots,Q;\ l=1,\ldots,L$ \\ \hline

6 &
$\displaystyle \log \kappa \sim \mathcal{N}(0,4)$ &
-- \\ \hline

\end{tabular}
\end{table}

여기서 $n$은 row node의 수, $p$는 column node의 수, $Q$는 row cluster의 수, $L$은 column cluster의 수를 의미한다. 또한
\[
\bar X = \frac{1}{np}\sum_{i=1}^{n}\sum_{j=1}^{p}X_{ij}
\]
는 전체 edge weight의 평균이며, prior strength는
\[
r = 0.1
\]
로 고정한다. 이 설정은 각 block parameter에 대하여 매우 약한 수준의 pseudo-information만을 제공한다.

\subsubsection{posterior distribution}

complete-data posterior는 다음과 같은 kernel을 가진다:
\begin{align*}
p(\mathbf{z},\mathbf{w},\pi,\rho,\Lambda,\kappa \mid X)
&\propto
\left[
\prod_{i=1}^{n}\prod_{j=1}^{p}
\frac{\lambda_{z_iw_j}^{\kappa}}{\Gamma(\kappa)}
X_{ij}^{\kappa-1}
\exp(-\lambda_{z_iw_j}X_{ij})
\right]
\left[
\prod_{i=1}^{n}\pi_{z_i}
\right]
\left[
\prod_{j=1}^{p}\rho_{w_j}
\right] \\
&\quad\times
\left[
\prod_{q=1}^{Q}\prod_{l=1}^{L}
\frac{(r\bar X)^{\kappa r}}{\Gamma(\kappa r)}
\lambda_{ql}^{\kappa r-1}
\exp(-r\bar X\,\lambda_{ql})
\right]
\operatorname{Dir}(\pi;1,\ldots,1)\,
\operatorname{Dir}(\rho;1,\ldots,1)\,
p(\kappa).
\end{align*}

표기 간소화를 위하여 다음의 sufficient statistics를 정의한다:
\[
n_q = \sum_{i=1}^{n}\mathbb{I}(z_i=q),
\qquad
m_l = \sum_{j=1}^{p}\mathbb{I}(w_j=l),
\]
\[
N_{ql}
=
\sum_{i=1}^{n}\sum_{j=1}^{p}\mathbb{I}(z_i=q,\; w_j=l),
\qquad
S_{ql}
=
\sum_{i=1}^{n}\sum_{j=1}^{p}\mathbb{I}(z_i=q,\; w_j=l)X_{ij}.
\]
완전 관측된 $n\times p$ bipartite matrix의 경우에는
\[
N_{ql} = n_q m_l
\]
가 성립한다.

\subsubsection{update rules}

사후 추론은 Gibbs sampler와 Metropolis--Hastings 알고리즘을 결합하여 수행한다. 구체적으로, $(\pi,\rho,\Lambda)$는 켤레 Gibbs 단계로 갱신하고, $(\mathbf{z},\mathbf{w})$는 categorical full conditional로부터 갱신하며, $\kappa$는 log scale random-walk MH 단계로 갱신한다.

\paragraph*{1. Update of $\pi$ (Gibbs step).}
Dirichlet--Categorical conjugacy에 의해
\[
\pi \mid \mathbf{z}
\sim
\operatorname{Dirichlet}(1+n_1,\ldots,1+n_Q).
\]

\paragraph*{2. Update of $\rho$ (Gibbs step).}
마찬가지로
\[
\rho \mid \mathbf{w}
\sim
\operatorname{Dirichlet}(1+m_1,\ldots,1+m_L).
\]

\paragraph*{3. Update of $\lambda_{ql}$ (Gibbs step).}
Gamma prior와 Gamma likelihood의 conjugacy에 의해, 각 block pair $(q,l)$에 대해
\[
\lambda_{ql}\mid X,\mathbf{z},\mathbf{w},\kappa
\sim
\operatorname{Gamma}
\left(
\kappa r + \kappa N_{ql},\;
r\bar X + S_{ql}
\right).
\]
즉, posterior shape는 prior contribution $\kappa r$와 observed block edge 수 $N_{ql}$의 합으로 결정되며, posterior rate는 prior contribution $r\bar X$와 block 내 관측된 edge weight의 총합 $S_{ql}$의 합으로 결정된다.

\paragraph*{4. Update of $z_i$ (categorical Gibbs step).}
각 $i=1,\ldots,n$에 대해, $z_i$의 full conditional은
\[
P(z_i=q \mid X,\mathbf{z}_{-i},\mathbf{w},\pi,\rho,\Lambda,\kappa)
\propto
\pi_q
\prod_{j=1}^{p}
f(X_{ij}\mid z_i=q,w_j,\kappa,\Lambda),
\qquad q=1,\ldots,Q.
\]
정규화에 필요하지 않은 상수항을 제거하면,
\[
P(z_i=q \mid \cdot)
\propto
\pi_q
\prod_{j=1}^{p}
\lambda_{q,w_j}^{\kappa}
\exp(-\lambda_{q,w_j}X_{ij}).
\]
따라서 계산의 안정성을 위해 log scale에서
\[
\log \omega_{iq}
=
\log \pi_q
+
\sum_{j=1}^{p}
\left[
\kappa \log \lambda_{q,w_j}
-
\lambda_{q,w_j}X_{ij}
\right]
\]
를 계산한 뒤,
\[
P(z_i=q\mid \cdot)=\frac{\exp(\log\omega_{iq})}{\sum_{q'=1}^{Q}\exp(\log\omega_{iq'})}
\]
로 정규화하여 categorical draw를 수행한다.

\paragraph*{5. Update of $w_j$ (categorical Gibbs step).}
각 $j=1,\ldots,p$에 대해,
\[
P(w_j=l \mid X,\mathbf{z},\mathbf{w}_{-j},\pi,\rho,\Lambda,\kappa)
\propto
\rho_l
\prod_{i=1}^{n}
f(X_{ij}\mid z_i,w_j=l,\kappa,\Lambda),
\qquad l=1,\ldots,L.
\]
상수항을 제거하면
\[
P(w_j=l \mid \cdot)
\propto
\rho_l
\prod_{i=1}^{n}
\lambda_{z_i,l}^{\kappa}
\exp(-\lambda_{z_i,l}X_{ij}),
\]
즉
\[
\log \upsilon_{jl}
=
\log \rho_l
+
\sum_{i=1}^{n}
\left[
\kappa \log \lambda_{z_i,l}
-
\lambda_{z_i,l}X_{ij}
\right]
\]
를 계산하고,
\[
P(w_j=l\mid \cdot)=\frac{\exp(\log\upsilon_{jl})}{\sum_{l'=1}^{L}\exp(\log\upsilon_{jl'})}
\]
로 정규화하여 categorical draw를 수행한다.

\paragraph*{6. Update of $\kappa$ (MH step on log scale).}
$\kappa>0$ 제약을 반영하기 위하여
\[
\xi = \log \kappa
\]
로 변환한 뒤, random-walk proposal
\[
\xi^{*} \sim \mathcal{N}(\xi,\sigma_{\text{prop},\kappa}^{2})
\]
를 사용한다. 이때
\[
\kappa^{*} = \exp(\xi^{*})
\]
이며, $\kappa$의 full conditional kernel은
\[
\tilde p(\kappa\mid \cdot)
\propto
p(\kappa)
\left[
\prod_{q=1}^{Q}\prod_{l=1}^{L}
p(\lambda_{ql}\mid \kappa)
\right]
\left[
\prod_{i=1}^{n}\prod_{j=1}^{p}
f(X_{ij}\mid z_i,w_j,\kappa,\Lambda)
\right]
\]
로 주어진다. 따라서 acceptance probability는
\[
\alpha_{\kappa}
=
\min\left\{
1,\;
\frac{\tilde p(\kappa^{*}\mid \cdot)\,\kappa^{*}}
{\tilde p(\kappa\mid \cdot)\,\kappa}
\right\}.
\]
여기서 $\kappa^{*}/\kappa$ 항은 log scale proposal에 따른 Jacobian correction에 해당한다.

\paragraph*{7. One full MCMC sweep.}
한 번의 전체 MCMC iteration은 다음 순서로 수행된다:
\[
\pi
\rightarrow
\rho
\rightarrow
\Lambda
\rightarrow
\mathbf{z}
\rightarrow
\mathbf{w}
\rightarrow
\kappa.
\]
실제로는 $\mathbf{z}$와 $\mathbf{w}$를 개별 성분별(component-wise)로 순차 업데이트한다.

\subsubsection{label switching post-processing}

본 모형은 row label과 column label의 순열에 대해 불변이므로, 저장된 MCMC 표본에 대하여 사후적 relabeling을 수행한다. 본 연구에서는 column-side clustering이 주요 해석 대상이므로, 먼저 각 저장 표본의 column membership $\mathbf{w}^{(t)}$를 참조 partition에 정렬한 뒤, 이에 맞추어 row membership $\mathbf{z}^{(t)}$를 정렬한다. 그 다음 block parameter matrix $\Lambda^{(t)}$에도 동일한 row/column permutation을 적용한다. 최종 posterior mean, credible interval, co-clustering probability, 그리고 block summary는 이러한 정렬을 마친 표본을 기준으로 계산한다.

\subsubsection{ICL-based selection of $(Q,L)$}

후보 모형 $(Q,L)$별로 MCMC를 수행한 뒤, complete-data log-likelihood에 BIC형 penalty를 부여한 ICL criterion을 계산한다:
\[
\operatorname{ICL}(Q,L)
=
\log p(X,\hat{\mathbf{z}},\hat{\mathbf{w}}\mid \hat{\Theta}_{Q,L},Q,L)
-
\frac{\nu_{Q,L}}{2}\log(np),
\]
where
\[
\nu_{Q,L}=(Q-1)+(L-1)+QL+1.
\]
여기서 마지막 $+1$은 global shape parameter $\kappa$에 해당한다. 최종 cluster 수는
\[
(\hat Q,\hat L)
=
\arg\max_{(Q,L)\in\mathcal{Q}\times\mathcal{L}}
\operatorname{ICL}(Q,L)
\]
로 선택한다. 본 연구의 목적상 column-side clustering이 더 중요하므로, 실제 탐색에서는 $L$에 대해 보다 넓거나 촘촘한 후보 집합을 둘 수 있다.

\subsubsection{tuning parameters}

MCMC iteration 수, burn-in 길이, thinning 간격, 그리고 MH proposal variance $\sigma_{\text{prop},\kappa}^{2}$는 모형 자체의 일부가 아니라 알고리즘 tuning quantity이므로 본 모형 기술에서는 고정하지 않는다. 이러한 값들은 실제 데이터의 크기와 mixing behavior를 고려하여 사용자가 결정한다.



\subsection{Post-LSIRM bipartite SBM using posterior distance samples}

본 절에서는 multilayered LSIRM으로부터 얻어진 item--respondent distance의 MCMC sample을
그대로 입력으로 사용하는 post-LSIRM bipartite SBM 절차를 제시한다.
multilayered LSIRM의 $s$번째 MCMC sample에서 응답자 $i$와 문항 $j$ 사이의 거리를
\[
D^{(s)}_{ij}
=
\left\| a_i^{(s)} - b_j^{(s)} \right\|,
\qquad
i=1,\ldots,n,\quad j=1,\ldots,p,\quad s=1,\ldots,m,
\]
라고 하자. 여기서 $n$은 respondent node의 수, $p$는 layer 구분 없이 결합한 전체 item node의 수,
$m$은 저장된 LSIRM MCMC sample의 수이다. 따라서 LSIRM posterior distance sample은
\[
\mathcal{D}
=
\left\{D^{(s)}\right\}_{s=1}^m,
\qquad
D^{(s)} \in \mathbb{R}_{+}^{n \times p}.
\]
각 원소 $D^{(s)}_{ij}$는 strictly positive distance-type edge weight로 해석되며, 값이 작을수록
respondent $i$와 item $j$가 latent space 상에서 더 가깝다는 것을 의미한다.

\subsubsection{Conditional SBM given each LSIRM distance sample}

본 연구에서는 LSIRM posterior uncertainty를 보존하기 위하여 하나의 posterior mean distance matrix를
사용하지 않고, 각 LSIRM MCMC sample $D^{(s)}$에 대해 별도의 bipartite SBM MCMC chain을 수행한다.
즉, $s$번째 LSIRM sample에 대해
\[
X^{(s)} = D^{(s)}
\]
로 두고, 다음의 bipartite SBM을 조건부로 적합한다.

row node와 column node의 cluster membership을 각각
\[
z_i^{(s)} \in \{1,\ldots,Q\},
\qquad
w_j^{(s)} \in \{1,\ldots,L\}
\]
로 둔다. 여기서 $Q$는 respondent-side cluster의 수이고, $L$은 item-side cluster의 수이다.
본 연구의 주요 관심은 item-side clustering이므로 최종 해석은 주로
\[
w^{(s)} = \left(w_1^{(s)},\ldots,w_p^{(s)}\right)
\]
를 중심으로 수행한다.

각 block pair $(q,l)$에 대해 distance edge는 Gamma likelihood를 따른다고 가정한다.
\[
X^{(s)}_{ij}
\mid
z_i^{(s)}=q,\,
w_j^{(s)}=l,\,
\kappa^{(s)},\,
\lambda_{ql}^{(s)}
\sim
\mathrm{Gamma}\left(\kappa^{(s)}, \lambda_{ql}^{(s)}\right),
\]
where $\kappa^{(s)} > 0$ is a global shape parameter and
$\lambda_{ql}^{(s)} > 0$ is a block-specific rate parameter. The corresponding density is
\[
f\left(X^{(s)}_{ij} \mid z_i^{(s)}=q,w_j^{(s)}=l,\kappa^{(s)},\lambda_{ql}^{(s)}\right)
=
\frac{\left(\lambda_{ql}^{(s)}\right)^{\kappa^{(s)}}}{\Gamma(\kappa^{(s)})}
\left(X^{(s)}_{ij}\right)^{\kappa^{(s)}-1}
\exp\left\{-\lambda_{ql}^{(s)}X^{(s)}_{ij}\right\}.
\]
Since
\[
\mathbb{E}\left(X^{(s)}_{ij} \mid z_i^{(s)}=q,w_j^{(s)}=l\right)
=
\frac{\kappa^{(s)}}{\lambda_{ql}^{(s)}},
\]
larger $\lambda_{ql}^{(s)}$ implies smaller average distance between respondent cluster $q$ and item cluster $l$.

\subsubsection{Prior specification}

For each LSIRM distance sample $s$, the prior distributions are specified as
\[
\pi^{(s)} \sim \mathrm{Dirichlet}(1,\ldots,1),
\qquad
\rho^{(s)} \sim \mathrm{Dirichlet}(1,\ldots,1),
\]
\[
z_i^{(s)} \mid \pi^{(s)} \sim \mathrm{Categorical}(\pi^{(s)}),
\qquad
i=1,\ldots,n,
\]
\[
w_j^{(s)} \mid \rho^{(s)} \sim \mathrm{Categorical}(\rho^{(s)}),
\qquad
j=1,\ldots,p.
\]
For each block-specific rate parameter,
\[
\lambda_{ql}^{(s)} \mid \kappa^{(s)}
\sim
\mathrm{Gamma}\left(\kappa^{(s)} r,\, r\bar X^{(s)}\right),
\qquad
q=1,\ldots,Q,\quad l=1,\ldots,L,
\]
where
\[
\bar X^{(s)}
=
\frac{1}{np}
\sum_{i=1}^n
\sum_{j=1}^p
X^{(s)}_{ij}.
\]
The prior strength is fixed at
\[
r = 0.1.
\]
Finally, the global shape parameter is assigned a weakly informative log-normal prior:
\[
\log \kappa^{(s)} \sim N(0,4).
\]

\subsubsection{Posterior distribution}

For each $s=1,\ldots,m$, the complete-data posterior is
\[
p\left(
z^{(s)},w^{(s)},\pi^{(s)},\rho^{(s)},\Lambda^{(s)},\kappa^{(s)}
\mid
X^{(s)}
\right)
\propto
p\left(X^{(s)} \mid z^{(s)},w^{(s)},\Lambda^{(s)},\kappa^{(s)}\right)
p\left(z^{(s)} \mid \pi^{(s)}\right)
p\left(w^{(s)} \mid \rho^{(s)}\right)
p\left(\pi^{(s)}\right)
p\left(\rho^{(s)}\right)
p\left(\Lambda^{(s)} \mid \kappa^{(s)}\right)
p\left(\kappa^{(s)}\right).
\]

Define the sufficient statistics
\[
n_q^{(s)}
=
\sum_{i=1}^n I\left(z_i^{(s)}=q\right),
\qquad
m_l^{(s)}
=
\sum_{j=1}^p I\left(w_j^{(s)}=l\right),
\]
\[
N_{ql}^{(s)}
=
\sum_{i=1}^n
\sum_{j=1}^p
I\left(z_i^{(s)}=q,w_j^{(s)}=l\right),
\]
\[
S_{ql}^{(s)}
=
\sum_{i=1}^n
\sum_{j=1}^p
I\left(z_i^{(s)}=q,w_j^{(s)}=l\right)
X^{(s)}_{ij}.
\]

\subsubsection{MCMC update rules for each LSIRM sample}

For each LSIRM distance sample $s$, one full MCMC sweep of the bipartite SBM is given by
\[
\pi^{(s)}
\rightarrow
\rho^{(s)}
\rightarrow
\Lambda^{(s)}
\rightarrow
z^{(s)}
\rightarrow
w^{(s)}
\rightarrow
\kappa^{(s)}.
\]

\paragraph{1. Update of $\pi^{(s)}$.}
By Dirichlet--Categorical conjugacy,
\[
\pi^{(s)} \mid z^{(s)}
\sim
\mathrm{Dirichlet}
\left(
1+n_1^{(s)},\ldots,1+n_Q^{(s)}
\right).
\]

\paragraph{2. Update of $\rho^{(s)}$.}
Similarly,
\[
\rho^{(s)} \mid w^{(s)}
\sim
\mathrm{Dirichlet}
\left(
1+m_1^{(s)},\ldots,1+m_L^{(s)}
\right).
\]

\paragraph{3. Update of $\lambda_{ql}^{(s)}$.}
For each block pair $(q,l)$,
\[
\lambda_{ql}^{(s)}
\mid
X^{(s)},z^{(s)},w^{(s)},\kappa^{(s)}
\sim
\mathrm{Gamma}
\left(
\kappa^{(s)}r+\kappa^{(s)}N_{ql}^{(s)},
\,
r\bar X^{(s)}+S_{ql}^{(s)}
\right).
\]

\paragraph{4. Update of $z_i^{(s)}$.}
For each respondent node $i=1,\ldots,n$,
\[
P\left(z_i^{(s)}=q \mid \cdot \right)
\propto
\pi_q^{(s)}
\prod_{j=1}^p
f\left(
X^{(s)}_{ij}
\mid
z_i^{(s)}=q,\,
w_j^{(s)},\,
\kappa^{(s)},\,
\Lambda^{(s)}
\right).
\]
Removing constants that do not depend on $q$, compute
\[
\log \omega_{iq}^{(s)}
=
\log \pi_q^{(s)}
+
\sum_{j=1}^p
\left[
\kappa^{(s)}
\log \lambda_{q,w_j^{(s)}}^{(s)}
-
\lambda_{q,w_j^{(s)}}^{(s)}
X^{(s)}_{ij}
\right].
\]
Then normalize:
\[
P\left(z_i^{(s)}=q \mid \cdot \right)
=
\frac{\exp\left(\log \omega_{iq}^{(s)}\right)}
{\sum_{q'=1}^Q \exp\left(\log \omega_{iq'}^{(s)}\right)}.
\]

\paragraph{5. Update of $w_j^{(s)}$.}
For each item node $j=1,\ldots,p$,
\[
P\left(w_j^{(s)}=l \mid \cdot \right)
\propto
\rho_l^{(s)}
\prod_{i=1}^n
f\left(
X^{(s)}_{ij}
\mid
z_i^{(s)},\,
w_j^{(s)}=l,\,
\kappa^{(s)},\,
\Lambda^{(s)}
\right).
\]
Equivalently,
\[
\log \upsilon_{jl}^{(s)}
=
\log \rho_l^{(s)}
+
\sum_{i=1}^n
\left[
\kappa^{(s)}
\log \lambda_{z_i^{(s)},l}^{(s)}
-
\lambda_{z_i^{(s)},l}^{(s)}
X^{(s)}_{ij}
\right].
\]
Then normalize:
\[
P\left(w_j^{(s)}=l \mid \cdot \right)
=
\frac{\exp\left(\log \upsilon_{jl}^{(s)}\right)}
{\sum_{l'=1}^L \exp\left(\log \upsilon_{jl'}^{(s)}\right)}.
\]

\paragraph{6. Update of $\kappa^{(s)}$.}
Let
\[
\xi^{(s)}=\log \kappa^{(s)}.
\]
Use a random-walk Metropolis--Hastings proposal
\[
\xi^{(s)\ast}
\sim
N\left(\xi^{(s)},\sigma^2_{\mathrm{prop},\kappa}\right),
\qquad
\kappa^{(s)\ast}=\exp\left(\xi^{(s)\ast}\right).
\]
The full conditional kernel of $\kappa^{(s)}$ is
\[
\widetilde p\left(\kappa^{(s)} \mid \cdot \right)
\propto
p\left(\kappa^{(s)}\right)
\prod_{q=1}^Q
\prod_{l=1}^L
p\left(\lambda_{ql}^{(s)} \mid \kappa^{(s)}\right)
\prod_{i=1}^n
\prod_{j=1}^p
f\left(
X^{(s)}_{ij}
\mid
z_i^{(s)},w_j^{(s)},\kappa^{(s)},\Lambda^{(s)}
\right).
\]
The proposal is accepted with probability
\[
\alpha_\kappa
=
\min
\left[
1,\,
\frac{
\widetilde p\left(\kappa^{(s)\ast}\mid\cdot\right)\kappa^{(s)\ast}
}{
\widetilde p\left(\kappa^{(s)}\mid\cdot\right)\kappa^{(s)}
}
\right],
\]
where the factor $\kappa^{(s)\ast}/\kappa^{(s)}$ is the Jacobian correction induced by the log-scale proposal.

\subsubsection{Pseudo-code}

\begin{algorithm}[H]
\caption{Post-LSIRM bipartite SBM using posterior distance samples}
\begin{algorithmic}[1]
\Require LSIRM posterior distance samples $\mathcal{D}=\{D^{(s)}\}_{s=1}^m$,
candidate cluster numbers $Q \in \mathcal{Q}$ and $L \in \mathcal{L}$,
SBM iteration number $T$, burn-in $B$, thinning interval $H$.
\For{$(Q,L) \in \mathcal{Q} \times \mathcal{L}$}
    \For{$s=1,\ldots,m$}
        \State Set $X^{(s)} \leftarrow D^{(s)}$.
        \State Compute $\bar X^{(s)} = (np)^{-1}\sum_{i=1}^n\sum_{j=1}^p X^{(s)}_{ij}$.
        \State Initialize $z^{(s,0)}$, $w^{(s,0)}$, $\pi^{(s,0)}$, $\rho^{(s,0)}$,
        $\Lambda^{(s,0)}$, and $\kappa^{(s,0)}$.
        \For{$t=1,\ldots,T$}
            \State Sample $\pi^{(s,t)} \mid z^{(s,t-1)}$ from its Dirichlet full conditional.
            \State Sample $\rho^{(s,t)} \mid w^{(s,t-1)}$ from its Dirichlet full conditional.
            \State Sample each $\lambda_{ql}^{(s,t)}$ from its Gamma full conditional.
            \For{$i=1,\ldots,n$}
                \State Sample $z_i^{(s,t)}$ from its categorical full conditional.
            \EndFor
            \For{$j=1,\ldots,p$}
                \State Sample $w_j^{(s,t)}$ from its categorical full conditional.
            \EndFor
            \State Update $\kappa^{(s,t)}$ by log-scale random-walk MH.
            \If{$t>B$ and $(t-B)$ is divisible by $H$}
                \State Store $\left(z^{(s,t)},w^{(s,t)},\pi^{(s,t)},\rho^{(s,t)},\Lambda^{(s,t)},\kappa^{(s,t)}\right)$.
            \EndIf
        \EndFor
    \EndFor
    \State Align stored labels by post-processing:
    first align item labels $w$, then respondent labels $z$, and permute $\Lambda$ accordingly.
    \State Compute posterior summaries and ICL$(Q,L)$.
\EndFor
\State Select
\[
(\widehat Q,\widehat L)
=
\arg\max_{(Q,L)\in\mathcal{Q}\times\mathcal{L}}
\mathrm{ICL}(Q,L).
\]
\State Report item-side posterior co-clustering probability
\[
\widehat P(w_j=w_{j'} \mid \mathcal{D})
=
\frac{1}{M_{\mathrm{save}}}
\sum_{\mathrm{saved}\, (s,t)}
I\left(w_j^{(s,t)}=w_{j'}^{(s,t)}\right),
\]
where $M_{\mathrm{save}}$ is the total number of saved post-LSIRM SBM samples.
\end{algorithmic}
\end{algorithm}

\subsubsection{ICL criterion}

For each candidate $(Q,L)$, define a representative partition
$(\widehat z,\widehat w)$ from the relabeled posterior samples.
The ICL criterion is computed as
\[
\mathrm{ICL}(Q,L)
=
\log p(\mathcal{D},\widehat z,\widehat w \mid \widehat\Theta_{Q,L},Q,L)
-
\frac{\nu_{Q,L}}{2}\log(mnp),
\]
where
\[
\widehat\Theta_{Q,L}
=
(\widehat\pi,\widehat\rho,\widehat\Lambda,\widehat\kappa),
\qquad
\nu_{Q,L}
=
(Q-1)+(L-1)+QL+1.
\]
Since $\mathcal{D}$ contains $m$ LSIRM posterior distance matrices, the completed log-likelihood is aggregated over
all LSIRM distance samples:
\[
\log p(\mathcal{D},\widehat z,\widehat w \mid \widehat\Theta_{Q,L},Q,L)
=
\sum_{s=1}^m
\left[
\sum_{i=1}^n \log \widehat\pi_{\widehat z_i^{(s)}}
+
\sum_{j=1}^p \log \widehat\rho_{\widehat w_j^{(s)}}
+
\sum_{i=1}^n
\sum_{j=1}^p
\log f\left(
D^{(s)}_{ij}
\mid
\widehat z_i^{(s)},\widehat w_j^{(s)},\widehat\kappa,\widehat\Lambda
\right)
\right].
\]
Finally,
\[
(\widehat Q,\widehat L)
=
\arg\max_{(Q,L)\in\mathcal{Q}\times\mathcal{L}}
\mathrm{ICL}(Q,L).
\]

\subsubsection{Remark}

이 절차는 multilayered LSIRM의 posterior distance sample을 단순히 평균내지 않고,
각 posterior draw를 하나의 가능한 distance network realization으로 취급한다.
따라서 최종 item clustering은 LSIRM latent position의 posterior uncertainty와
bipartite SBM clustering uncertainty를 동시에 반영한다.
\end{document}